% =============================================================================
% Article C (JEPA-Drift) -- Appendix A, proofs of Section 3, v0 (2026-10-15).
% =============================================================================
\section{Proofs of Section~\ref{sec:theory}}\label{app:proofs}

\paragraph{Proposition~\ref{prop:hierarchy}.}
(i)~The map $(a,S)\mapsto a^{\top}S^{-1}a$ is jointly convex on $\mathbb{R}^N\times\mathbb{S}^N_{++}$ \citep[\S3.1.7]{boyd_vandenberghe_2004}. Jensen's inequality over the law of $(\Delta\mu_t,\Sigma_t)$ gives the first inequality. For a fixed $w$, the Cauchy--Schwarz inequality in the $\bar\Sigma$ inner product gives $(w^{\top}\overline{\Delta\mu})^2/(2w^{\top}\bar\Sigma w)\le\tfrac12\overline{\Delta\mu}^{\top}\bar\Sigma^{-1}\overline{\Delta\mu}$; the left side is the Gaussian quasi-divergence of the projected stream.
(ii)~Let $f_t$ be the conditional density of $x_t$ and $\varsigma=\nabla_x\log f_t(x_t)$ its location score. Integration by parts gives $\mathbb{E}[\varsigma\,(x_t-\mu_t)^{\top}\mid\mathcal{F}_{t-1}]=-I_N$, and $\mathcal{J}_t=\mathbb{E}[\varsigma\varsigma^{\top}\mid\mathcal{F}_{t-1}]$. The second-moment matrix of $(\varsigma,x_t-\mu_t)$ is positive semidefinite; its Schur complement is $\mathcal{J}_t-\Sigma_t^{-1}\succeq0$, with equality if and only if $\varsigma$ is affine in $x_t$, that is, if $f_t$ is Gaussian.
(iii)~Under the usual regularity of location families, $\mathrm{KL}\big(f_t(\cdot-\Delta)\,\|\,f_t\big)=\tfrac12\Delta^{\top}\mathcal{J}_t\Delta+o(\|\Delta\|^2)$ \citep{kullback_1959}. The chain rule adds these conditional divergences over $t$, and the ergodic theorem turns their average into $\mathcal{I}_x$.
(iv)~For every $n$, $z_{1:n}$ is a measurable function of the monitored data $(x_{1:n},C_{1:n})$, and the data-processing inequality \citep{cover_thomas_2006} gives $\mathrm{KL}(P_1^{z_{1:n}}\|P_0^{z_{1:n}})\le\mathrm{KL}(P_1^{x_{1:n},C_{1:n}}\|P_0^{x_{1:n},C_{1:n}})$. Dividing by $n$ and letting $n\to\infty$ yields $\mathcal{I}_z\le\mathcal{I}_x$. \hfill$\square$

\paragraph{Lemma~\ref{lem:budget}.}
Write $\lambda_t=e^{\ell_t}$, $\Lambda_n=\prod_{t\le n}\lambda_t$ and $T=T_\gamma$. \emph{Step~1.} The ratio $R_n=E_n/\Lambda_n$ is a nonnegative $P_1$-supermartingale. Given $\mathcal{F}_{n-1}$, $dP_1=\lambda_n\,dP_0$ on $\mathcal{F}_n$, so
$\mathbb{E}_1[R_n\mid\mathcal{F}_{n-1}]=\Lambda_{n-1}^{-1}\,\mathbb{E}_0\big[E_n\mathbf 1\{\lambda_n>0\}\mid\mathcal{F}_{n-1}\big]\le E_{n-1}/\Lambda_{n-1}=R_{n-1}.$
\emph{Step~2.} Since $\mathbb{E}_1[T]<\infty$, $T$ is finite almost surely. Optional stopping for nonnegative supermartingales gives $\mathbb{E}_1[R_T]\le R_0\le1$, and Jensen's inequality gives $\mathbb{E}_1[\log R_T]\le0$.
\emph{Step~3.} The increments $d_t=\ell_t-I_t$ are martingale differences with $\mathbb{E}_1[|d_t|\mid\mathcal{F}_{t-1}]\le2c$, where $c=\sup_t\mathbb{E}_1[|\ell_t|\mid\mathcal{F}_{t-1}]$. Since $\mathbb{E}_1\big[\sum_{t\le T}|d_t|\big]\le2c\,\mathbb{E}_1[T]<\infty$, dominated convergence gives Wald's identity $\mathbb{E}_1[\log\Lambda_T]=\mathbb{E}_1\big[\sum_{t\le T}I_t\big]$.
\emph{Step~4.} On $\{T<\infty\}$, $\log E_T\ge\ln\gamma$ and $\log E_T=\log R_T+\log\Lambda_T$. Step~3 makes $\mathbb{E}_1[\log\Lambda_T]$ finite, so $\mathbb{E}_1[\log R_T]$ is finite as well, and
$\ln\gamma\le\mathbb{E}_1[\log E_T]\le\mathbb{E}_1[\log\Lambda_T]=\mathbb{E}_1\big[\sum_{t\le T}I_t\big]$.
For $E=\Lambda$, $R\equiv1$ and the only slack is $\mathbb{E}_1[\log\Lambda_T]-\ln\gamma$, the expected overshoot. \hfill$\square$

\paragraph{Proposition~\ref{prop:floor}.}
(a)~The CUSUM stopping time equals $\min_j\{j-1+N_j\}$, where $N_j$ starts at $j$ and stops when its likelihood ratio reaches $\gamma$. Ville's inequality gives $P_0(N_j<\infty)\le1/\gamma$, and \citet[Thm.~2]{lorden_1971} yields $\mathrm{ARL}_0\ge\gamma$. Under $P_1$, $\ell_t=\theta^{\top}z_t-\|\theta\|^2/2$ is Gaussian with mean $I_z$ and variance $2I_z$, and the CUSUM is the walk $\sum\ell_t$ reflected at zero; Lorden's worst case is the restart from zero. A reflected Brownian motion with drift $\mu$ and variance $\sigma^2$, started at zero, reaches $h$ after an expected time $\frac{\sigma^2}{2\mu^2}\big(e^{-2\mu h/\sigma^2}-1+2\mu h/\sigma^2\big)$ \citep{taylor_1975}. With $\mu=I_z$ and $\sigma^2=2I_z$ this equals $(h-1+e^{-h})/I_z$. The corrected diffusion approximation replaces $h$ by $h+2\rho\sigma_\ell$, where $\sigma_\ell=\|\theta\|$ is the standard deviation of $\ell_t$ and $\rho=0.5826$ \citep[Ch.~X]{siegmund_1985}.
(b)~Optimality is the theorem of \citet{moustakides_1986}. Under $P_0$ the walk has drift $-I_z$, and the same formula gives $I_z\,\mathrm{ARL}_0=e^{h}-h-1$; calibration to $\mathrm{ARL}_0=\gamma$ yields $e^{h}-h-1=\gamma I_z$ and $I_z\,\mathrm{ADD}=\Phi(\gamma I_z)$. As $I_z\to0$ with $\gamma I_z$ fixed, the scaled hitting times of the reflected walk converge to those of the reflected Brownian motion by the invariance principle, with uniform integrability of the hitting times \citep[Ch.~X]{siegmund_1985}. For large $g$, $h_g=\ln(g+h_g+1)=\ln g+O(g^{-1}\ln g)$ and $e^{-h_g}=O(g^{-1})$, so $\Phi(g)=\ln g-1+O(g^{-1}\ln g)$. For small $g$, $g=h^2/2+h^3/6+h^4/24+O(h^5)$ and $\Phi=h^2/2-h^3/6+h^4/24+O(h^5)$, hence $\Phi=g-h^3/3+O(h^5)$; with $h=\sqrt{2g}\,\big(1+O(\sqrt g)\big)$ this is $g-(2g)^{3/2}/3+O(g^2)$.
(c)~Part (ii) of Proposition~\ref{prop:hierarchy}, applied to the latent stream, gives the local divergence $\tfrac12\Delta\mu_z^{\top}\mathcal{J}_z\Delta\mu_z\ge\tfrac12\Delta\mu_z^{\top}\Sigma_z^{-1}\Delta\mu_z$, with equality in the Gaussian case. By Cauchy--Schwarz, $\max_w(w^{\top}\Delta\mu_z)^2/(w^{\top}\Sigma_zw)=\Delta\mu_z^{\top}\Sigma_z^{-1}\Delta\mu_z$, attained at $w^\ast=\Sigma_z^{-1}\Delta\mu_z$: the squared daily Sharpe ratio of the drift-tangent portfolio. Annualizing multiplies it by 252, so $I_z\ge\mathrm{SR}_z^2/504$. Proposition~\ref{prop:hierarchy} gives $I_z\le\mathcal{I}_x$, and $\mathcal{I}_x=\mathrm{SR}_x^2/504$ when $x_t$ is conditionally Gaussian. \hfill$\square$

\paragraph{Corollary~\ref{cor:factor}.}
The Woodbury identity gives $\Sigma^{-1}=\Psi^{-1}-\Psi^{-1}B(\Sigma_F^{-1}+K)^{-1}B^{\top}\Psi^{-1}$, hence $\Sigma^{-1}B=\Psi^{-1}B(\Sigma_F^{-1}+K)^{-1}\Sigma_F^{-1}$. By construction $B^{\top}\Psi^{-1}\Delta\alpha_\perp=0$, so $\Delta\alpha_\perp^{\top}\Sigma^{-1}B=0$ and $\Delta\alpha_\perp^{\top}\Sigma^{-1}\Delta\alpha_\perp=\Delta\alpha_\perp^{\top}\Psi^{-1}\Delta\alpha_\perp$. Writing $\Delta\mu=B\Delta\tilde\lambda+\Delta\alpha_\perp$ leaves the factor term $\Delta\tilde\lambda^{\top}B^{\top}\Sigma^{-1}B\,\Delta\tilde\lambda$, and $B^{\top}\Sigma^{-1}B=K(I+\Sigma_FK)^{-1}=(\Sigma_F+K^{-1})^{-1}\preceq\Sigma_F^{-1}$. With one factor, $B=\mathbf 1$, $\Sigma_F=\rho$, $\Psi=(1-\rho)I_N$ and $K=N/(1-\rho)$, the factor term is $\Delta m^2/\big(\rho+(1-\rho)/N\big)=\Delta m^2N/\big(1+(N-1)\rho\big)$, against $\Delta m^2$ for one asset of unit variance. \hfill$\square$

\paragraph{Proposition~\ref{prop:tax}.}
(a)~For fixed $u$, $M_n(u)=\exp\{u^{\top}S_n-n\|u\|^2/2\}$ is a $P_0$-supermartingale when $\mathbb{E}_0[e^{u^{\top}z_t}\mid\mathcal{F}_{t-1}]\le e^{\|u\|^2/2}$. Gaussian coordinates meet the bound with equality; independent Rademacher coordinates meet it because $\cosh v\le e^{v^2/2}$. Tonelli's theorem carries the supermartingale property to the mixture \citep{robbins_1970,howard_2020}, and the Gaussian integral gives $L_n=(1+ns^2)^{-D/2}\exp\{s^2\|S_n\|^2/(2(1+ns^2))\}$.
(b)~Under $P_1$, $\mathbb{E}_1\|S_n\|^2=n^2\|\theta\|^2+nD=2n^2I+nD$. With $s^2=2I/D$ and $x=nI$, $\mathbb{E}_1[\log L_n]=\frac{2x/D}{1+2x/D}\big(x+\tfrac D2\big)-\tfrac D2\ln(1+2x/D)=x-\tfrac D2\ln(1+2x/D)$.
(c)~For $D\gg x$, $\tfrac D2\ln(1+2x/D)=x-x^2/D+O(x^3/D^2)$, so $x^\ast\simeq\sqrt{D\ln\gamma}$. The fluid approximation neglects the overshoot and the fluctuations of $\|S_n\|^2$; Section~\ref{sec:exp-floors} measures its error.
(d)~An average of $K$ test supermartingales is a test supermartingale whatever their dependence, and it exceeds $\gamma$ as soon as one of them exceeds $K\gamma$. Along the unit atom $a_k$, the likelihood-ratio e-process with matched magnitude grows at rate $(a_k^{\top}\theta)^2/2=I\cos^2\vartheta$, so it reaches $K\gamma$ after $(\ln\gamma+\ln K)/(I\cos^2\vartheta)$ steps in the fluid approximation. \hfill$\square$

\paragraph{Corollary~\ref{cor:scaling}.}
With $D$ fixed, $M(D,\gamma)$ does not depend on $N$ and the fluid delay scales as $1/I(N)$. With $D\gg\ln\gamma$, Proposition~\ref{prop:tax}(c) gives a delay proportional to $\sqrt{D\ln\gamma}/I(N)$. The single-asset delay does not depend on $N$, so the gain scales as $N^{a}$ and $N^{a-b/2}$ respectively. A market drift has $a=0$ by Corollary~\ref{cor:factor}. \hfill$\square$

\paragraph{Lemma~\ref{lem:odd}.}
Conditionally on $\mathcal{F}_{t-1}$, $g(x)=u_t^{\top}\phi(x;C_t)$ is odd, so central symmetry gives $P_0(g(x_t)>0\mid\mathcal{F}_{t-1})=P_0(g(-x_t)>0\mid\mathcal{F}_{t-1})=P_0(g(x_t)<0\mid\mathcal{F}_{t-1})$. The two probabilities sum to one, hence equal $1/2$. Since $b_{1:t-1}$ is $\mathcal{F}_{t-1}$-measurable, $P_0(b_{1:t})=\tfrac12P_0(b_{1:t-1})$, and induction gives the i.i.d.\ Bernoulli$(1/2)$ law. \hfill$\square$

\paragraph{Corollary~\ref{cor:orth}.}
(i)~Lemma~\ref{lem:odd} applies before and after any change that preserves central symmetry, so the sign stream keeps the i.i.d.\ Bernoulli$(1/2)$ law. (ii)~If $s$ is even, $\nabla s$ is odd and $\mathbb{E}[\nabla s(x)]=0$ under central symmetry; a second-order Taylor expansion gives $\mathbb{E}[s(x+\delta)]-\mathbb{E}[s(x)]=\tfrac12\delta^{\top}\mathbb{E}[\nabla^2s(x)]\delta+o(\|\delta\|^2)$. \hfill$\square$

\paragraph{Proposition~\ref{prop:absorption}.}
(i)~Suppose that $T=T_\gamma<\infty$ $P_1$-almost surely. The martingale $\sum_{t\le n}(\ell_t-I_t)$ has conditional variances $2I_t$ that sum to $2B<\infty$; it is bounded in $L^2$, hence uniformly integrable, and optional stopping at $T$ gives $\mathbb{E}_1[\log\Lambda_T]=\mathbb{E}_1[\sum_{t\le T}I_t]\le B$. Steps~1 and~2 of the proof of Lemma~\ref{lem:budget} give $\mathbb{E}_1[\log R_T]\le0$. Then $\ln\gamma\le\mathbb{E}_1[\log E_T]\le B<\ln\gamma$, a contradiction.
(ii)~In continuous time the information accumulated after $t$ days is $\int_0^tI_ze^{-2s/\tau_a}ds=\tfrac12I_z\tau_a\big(1-e^{-2t/\tau_a}\big)$. Equating it with $\ln\gamma=I_z\tau_\gamma$ gives $t=-\tfrac{\tau_a}{2}\ln(1-2\tau_\gamma/\tau_a)$, finite if and only if $\tau_a>2\tau_\gamma$. \hfill$\square$
